
\documentclass[12pt,aspectratio=169,t,fleqn]{beamer}

\usepackage[fontset=none,autoindent=true]{ctex}
\IfFontExistsTF{Songti SC}{%
	\setCJKmainfont{Songti SC}[AutoFakeSlant=0.2]%
	\setCJKsansfont{Songti SC}[AutoFakeSlant=0.2]%
	\setCJKmonofont{Songti SC}%
}{%
	\setCJKmainfont{FandolSong-Regular}[Extension=.otf,BoldFont=FandolSong-Bold,AutoFakeSlant=0.2]%
	\setCJKsansfont{FandolSong-Regular}[Extension=.otf,BoldFont=FandolSong-Bold,AutoFakeSlant=0.2]%
	\setCJKmonofont{FandolFang-Regular}[Extension=.otf,BoldFont=FandolFang-Regular]%
}
\usepackage{hyperref}
\setmainfont{XCharter} % Nicer fonts — the charter package is ignored under XeLaTeX, so use fontspec instead
% 等宽字体：优先 Maple Mono NF（静态字族，真实粗/斜体），其次可变版，未安装则回退默认
\IfFontExistsTF{Maple Mono NF}{%
	\setmonofont{Maple Mono NF}[Scale=MatchLowercase,Ligatures=TeX]%
}{%
	\IfFontExistsTF{Maple Mono}{%
		\setmonofont{Maple Mono}[Scale=MatchLowercase,Ligatures=TeX]%
	}{}%
}
\usefonttheme{professionalfonts}%**强烈建议加上**，关闭 beamer 自带字体替换，防止字体冲突，数学公式字体正常渲染
\mode<presentation>
{
  \usetheme{Warsaw}

  \setbeamercovered{invisible}
 %   含义：**大括号内的代码，只在 presentation（幻灯片）模式下生效；
 %   如果切换成 article 讲义模式，这一段直接忽略不执行。**
 
% 控制**被覆盖（还没出现）的条目**样式： 
% `transparent`：未出现的条目变浅、半透明（灰色淡化，占位还在页面上）
% 可选：`invisible`（默认，没出现的内容直接隐藏，页面空白）
% 对应你之前的 `<+->` 逐条动画：
% `transparent`：后面还没出来的 \item 是淡淡的灰色，提前看到占位
% `invisible`：没到显示时机的 item 完全看不见
}


\title[2.1等式性质与不等式性质] % (optional, use only with long paper titles)
{2.1等式性质与不等式性质}

\subtitle
{等式性质与不等式性质}

\author[好好学习，努力向上] % (optional, use only with lots of authors)
{
%	金永权\inst{1} \and 某某某\inst{2}
	}

\institute[夏邑一高] % (optional, but mostly needed)
{
%  \inst{1}%
%  高一年级
%  22 班
%  \and
%  \inst{2}%
% 高一年级
% 24班
 }
% - Use the \inst command only if there are several affiliations.
% - Keep it simple, no one is interested in your street address.

\date[2026.9.19] % (optional, should be abbreviation of conference name)
{2026.9.19}
% - Either use conference name or its abbreviation.
% - Not really informative to the audience, more for people (including
%   yourself) who are reading the slides online

\subject{数学科}
% This is only inserted into the PDF information catalog. Can be left out. 


\AtBeginSubsection[]
{
  \begin{frame}<beamer>{主要内容}
    \tableofcontents[sectionstyle=show/hide,subsectionstyle=show/shaded]
  \end{frame}
}
%	作用：**每次开启一个 `\subsection` 小节时，自动执行大括号里面的代码**。
%	`[]` 里面可以放 overlay 规范，空括号代表默认。
%	含义：**这个目录帧只在 presentation（幻灯片模式）编译出来；
%	如果切换成 article 讲义模式，这一帧直接忽略，不会生成**。
%	等价写法：`\mode<presentation>{\begin{frame}...\end{frame}}`
%	目录命令的选项：

%	- `currentsection`：**只高亮当前所在的大 section，其他 section 变灰**
%	- `currentsubsection`：**高亮当前 subsection，同 section 下其他小节变浅**
%\beamerdefaultoverlayspecification{<+->}


\begin{document}

\begin{frame}
  \titlepage
\end{frame}

\begin{frame}{主要内容}
  \tableofcontents
  % You might wish to add the option [pausesections]
\end{frame}


\section{等式的基本性质}

\subsection{等式的基本性质}

\begin{frame}{等式的基本性质}
  % - A title should summarize the slide in an understandable fashion
  %   for anyone how does not follow everything on the slide itself.

  \begin{enumerate}
  \item    如果 $a=b$，那么$b=a$；
  \item    如果 $a=b$，$b=c$，那么 $a=c$；
  \item  	如果$a=b$，那么 $a\pm c = b \pm c$；
  \item  	如果 $a=b$，那么 $ac=bc$；
  \item 	如果 $a=b,c\neq 0$，那么$\frac{a}{c} = \frac{b}{c}$。
  \end{enumerate}
  
\end{frame}

\section{不等式的基本性质}
\subsection{不等式的基本性质}

\begin{frame}{不等式的基本性质}

  类比等式的性质1，2，可以猜想不等式有如下性质：
  \begin{enumerate}[<+->]
	\item 如果$a>b$，那么$b<a$；如果$b<a$，那么$a>b$。即
	 		 \[a>b \iff b<a\]
	    
	\item 如果$a>b$，$b>c$，那么$a>c$，即
	      \[ a>b,b>c \Rightarrow a>c\]

  \end{enumerate}
\end{frame}

\begin{frame}{不等式的基本性质}
类比等式的性质 3\textasciitilde 5，可以猜想不等式还有如下性质：
	\begin{enumerate}[<+->]\addtocounter{enumi}{2}
		
	\item  如果$a>b$，那么$a+c>b+c$；	
	\item 如果 $a>b,\ c>0$，那么 $ac>bc$；如果 $a>b,\ c<0$，那么 $ac<bc$.
	
	\item  如果 $a>b,\ c>d$，那么 $a+c>b+d$.

	\end{enumerate}
\end{frame}

\begin{frame}{不等式的基本性质}
	利用性质 4 和性质 2 可以推出：
	\begin{enumerate}[<+->]\addtocounter{enumi}{5}
		
		\item 如果 $a>b>0$，$c>d>0$，那么 $ac>bd$。
		
		\item 如果 $a>b>0$，那么 $a^n>b^n\ (n\in\mathbf{N},\ n\ge 2)$。
	\end{enumerate}
\end{frame}

\section{例题分析}
\begin{frame}{}
	\begin{example}
		已知 $a>b>0$，$c<0$，求证 $\dfrac{c}{a}>\dfrac{c}{b}$。
	\end{example}
	
	\pause

	证明：	因为 $a>b>0$，所以 $ab>0$，$\dfrac{1}{ab}>0$。
		
		于是
		\[
		a\cdot \frac{1}{ab}>b\cdot \frac{1}{ab},
		\]
		即
		\[
		\frac{1}{b}>\frac{1}{a}.
		\]
		
		由 $c<0$，得 $\dfrac{c}{a}>\dfrac{c}{b}$。

\end{frame}
\begin{frame}{布置作业}
	P42 练习: 1、2
	
	P42 习题2.1: 5，6，7，8	
\end{frame}


\end{document}


